\section{材料与方法}
\label{sec:methods}

\subsection{传播、报告与两类估计目标}
\label{sec:bg-model}

考虑以 $t$ 为索引的一代传播，其中 $n$ 名感染者作为原发病例。假设给定瞬时再生数 $R_t$ 时，原发病例 $i$ 引起的继发感染数 $X_{i,t}$ 服从负二项分布，以刻画传播中的个体异质性\citep{lloydsmith2005superspreading}。后代分布为
\begin{equation}
  X_{i,t}\mid R_t \sim \NB\!\left(k,\ \frac{k}{k+R_t}\right),\qquad % qa-allow
  \E(X_{i,t}\mid R_t)=R_t,\qquad \Var(X_{i,t}\mid R_t)=R_t+\frac{R_t{}^2}{k}, % qa-allow
  \label{eq:offspring}
\end{equation}
其中 $k>0$ 为离散参数，$k$ 越小，异质性越强；$k\to\infty$ 时退化为 Poisson 后代分布。每例感染随后以概率 $\rho\in(0,1]$ 独立报告。成功概率相同的独立负二项变量之和仍服从负二项分布，二项稀疏化也保持这一分布族，因此报告病例数 $C_t$ 满足
\begin{equation}
  C_t\mid n,R_t \sim \NB\!\left(nk,\ \frac{k}{k+\rho R_t}\right),\qquad % qa-allow
  \E(C_t\mid n,R_t)=\rho nR_t,\qquad % qa-allow
  \Var(C_t\mid n,R_t)=\rho nR_t+\frac{\rho^2nR_t{}^2}{k},
  \label{eq:observed}
\end{equation}
其中 $m=\rho n$ 为 $R_t=1$ 时每个时期的期望报告病例数，下文称 $m$ 为病例量。

模型建立在四项假设之上。每例感染以不随时间变化的概率 $\rho$ 独立报告（H1）。原发病例按式~(\ref{eq:offspring}) 独立同分布地传播（H2）。离散参数 $k$ 与 $R_t$ 和 $n$ 无关（H3）。原发病例数 $n$ 由设计固定，且与 $R_t$ 相互独立（H4）。当 $n$ 由接触者追踪确定时，假设 H4 成立；常规监测中 $n$ 本身是早期传播的结果，第~\ref{sec:meth-loss}~节的更新方程形式放宽了这一假设。

传播潜力随天气、学期、节假日和人群行为在时期之间变化\citep{dalziel2018urbanization}，其边际矩记为 $\E(R_t)=\Rbar$ 与 $\Var(R_t)=v_R$。这类变异同时作用于同一时期内的每条传播链，对应种群生态学中的环境随机性；式~(\ref{eq:observed}) 中的后代变异与报告变异则对应种群统计随机性\citep{engen1998demographic, lande2003stochastic}。

尺度化病例数 $\Rhat_t=C_t/m$ 可服务于两个不同的估计目标。第一个目标是当期值 $R_t$，用于态势感知，也用于判断当前 $R_t$ 是否大于一。第二个目标是跨期均值 $\Rbar$，用于描述一个流行季内的典型传播水平；当 $R_t$ 在各时期间相互独立时，$\Rbar$ 也是下一时期取值的最优一步预测。第一个目标的误差为 $\Rhat_t-R_t$，第二个目标的误差可写为
\begin{equation}
  \Rhat_t-\Rbar=(\Rhat_t-R_t)+(R_t-\Rbar),
  \label{eq:error-split}
\end{equation}
其中第二项与 $m$ 无关。给定 $R_t$ 时，$\Rhat_t$ 是 $R_t$ 的无偏估计，其方差 $R_t/m+\rho R_t{}^2/(km)$ 随 $m$ 增大趋于零，并达到 $R_t$ 的 Cram\'er--Rao 下界（附录~\ref{app:target1}）。

\subsection{误差方差与平衡规模}
\label{sec:bg-variance}

对第二个目标，将全方差公式应用于式~(\ref{eq:error-split}) 即可得到误差方差。

\begin{proposition}[跨期均值的误差方差]
\label{prop:average}
在 H1--H4 下，$\Rhat_t=C_t/(\rho n)$ 围绕 $\Rbar$ 的方差为
\begin{equation}
  V(n,\rho)=v_R+\frac{\Rbar}{\rho n}+\frac{\Rbar^2+v_R}{kn}=v_R+\frac{A}{m},
  \qquad A=\Rbar+\frac{\rho(\Rbar^2+v_R)}{k}, % qa-allow
  \label{eq:var-average}
\end{equation}
其中三项依次为环境贡献、报告贡献与后代贡献。由式~(\ref{eq:var-average}) 可得以下四点。
\begin{enumerate}[label=(\roman*),leftmargin=2em]
  \item 对固定的 $\rho$，当 $n\to\infty$ 时 $V\to v_R$。
  \item 对固定的 $n$，当 $\rho\to1$ 时 $V\to v_R+\Rbar/n+(\Rbar^2+v_R)/(kn)>v_R$。
  \item 固定 $\rho$、$k$ 与 $R_t$ 的分布，经由 $n$ 改变 $m=\rho n$ 时，在平衡规模 $\mx=A/v_R$ 处抽样贡献与环境贡献相等；$m$ 加倍使 $V$ 按比例 $1/[2(1+m/\mx)]$ 减小，在 $m=\mx$ 处为四分之一。
  \item 在对 $R_t$ 条件无偏的估计量中，任何估计量对 $\Rbar$ 的均方误差均不小于 $V$。
\end{enumerate}
\end{proposition}

第（iii）项中的 $A$ 依赖于 $\rho$，因此平衡规模和加倍收益只对固定 $\rho$ 的路径成立。若固定 $n$ 而把 $\rho$ 提高到 $2\rho$（要求 $\rho\le1/2$），只有报告项减半，方差的相对降幅为 $\Rbar/[2\rho nV(n,\rho)]$，后代项 $(\Rbar^2+v_R)/(kn)$ 不变；这一路径的可达下限为第（ii）项给出的取值。两条路径的对数边际分别为 $-\partial V/\partial\log n=\Rbar/(\rho n)+(\Rbar^2+v_R)/(kn)=A/m$ 和 $-\partial V/\partial\log\rho=\Rbar/(\rho n)=\Rbar/m$，二者之比为 $A/\Rbar$（附录~\ref{app:prop1}）。

平衡规模的作用相当于生态学模型中的临界种群规模；在临界种群规模处，种群增长方差的种群统计贡献与环境贡献相等\citep{lande2003stochastic}。平衡规模是比较两类贡献的参照点，本身不界定最低运行规模。第（iv）项仅适用于条件无偏估计量。将 $\Rhat_t$ 向先验均值收缩的估计量对 $\Rbar$ 的误差可能更小（附录~\ref{app:shrinkage}）；对 $T$ 个独立时期取平均，可将环境项降至 $v_R/T$（命题~\ref{prop:allocation}）。极限 $v_R$ 还假设环境变异在时期之间和地点之间相互独立。若 $R_t$ 服从自相关系数为 $\phi$ 的平稳一阶自回归过程，一步预测误差中的环境部分降至 $(1-\phi^2)v_R$。若 $J$ 个辖区的环境方差均为 $v_R$，两两相关系数的均值为 $\bar{c}$，则这些辖区平均值的环境方差为
\begin{equation}
  v_R\,\bar{c}+\frac{v_R(1-\bar{c})}{J},
  \label{eq:pooling}
\end{equation}
即汇总仅在变异不同步（$\bar{c}<1$）的范围内降低这一极限。

\subsection{报告异质性与监测预算分配}
\label{sec:bg-reporting}

报告完整度很少保持恒定。假设 H1 可替换为如下假设：报告概率 $p_t$ 在时期之间变化，且与 $R_t$、$n$ 相互独立，其均值为 $\rho$，相对波动为 $\tau=\mathrm{SD}(p_t)/\rho$；给定 $R_t$ 与 $p_t$ 时，各时期病例数相互独立（H5）。

\begin{proposition}[报告异质性]
\label{prop:report-noise}
在 H2--H5 下，$\Rhat_t=C_t/m$ 围绕 $\Rbar$ 的误差方差可写为
\begin{equation}
  V_H(m)=v_R+v_H+\frac{A_H}{m},\qquad % qa-allow
  v_H=\tau^2(\Rbar^2+v_R),\qquad A_H=\Rbar+\frac{\rho(1+\tau^2)(\Rbar^2+v_R)}{k}, % qa-allow
  \label{eq:vH-AH}
\end{equation}
其中 $v_H$ 为报告异质性引入的附加方差。极限方差升至 $v_R+v_H$，平衡规模降至 $m_{\times,H}=A_H/(v_R+v_H)<\mx$。仅凭病例数的均值与方差无法区分 $v_R$ 与 $v_H$，这两个矩只能识别 $(1+\tau^2)(\Rbar^2+v_R)$。
\end{proposition}

命题~\ref{prop:report-noise} 构成实证部分的识别边界：若不单独测量报告过程，仅由病例数估计的截距同时包含传播变异与报告变异。实际增长统计量的分子与分母来自同一报告序列，报告冲击还会在相邻周之间产生负的协方差，并通过分母进入均值；附录~\ref{app:random-denominator} 给出随机分子分母下的二阶展开及其对增长比与更新型比值的含义。

平衡规模还决定固定预算的最优分配方式。设共观测 $T$ 个时期，每个时期的病例量为 $m$，每个时期的成本为 $c_0$，每例报告病例的成本为 $c_1$，总预算为 $\mathcal{B}=T(c_0+c_1m)$。

\begin{proposition}[固定预算的分配]
\label{prop:allocation}
在 H1--H4 下固定 $\rho$，经由 $n$ 改变 $m$。若 $R_t$ 在各时期间相互独立，则 $T$ 个时期内 $\Rhat_t$ 均值的方差为
\begin{equation}
  \frac{v_R}{T}+\frac{A}{Tm},
  \label{eq:alloc-indep}
\end{equation}
当 $c_0,c_1,v_R,A>0$ 且连续内点可行时，预算约束下上述方差的唯一最小值点为
\begin{equation}
  m^*=\sqrt{\frac{c_0}{c_1}\,\mx},
  \label{eq:optimal-m}
\end{equation}
其中 $T^*=\mathcal{B}/(c_0+c_1m^*)$，并要求 $T^*\ge1$。若 $R_t$ 服从自相关系数为 $\phi$ 的平稳一阶自回归过程，抽样项保持不变，环境项变为 $v_RS(\phi,T)/T^2$，其中 $S(\phi,T)$ 为
\begin{equation}
  S(\phi,T)=\frac{T(1+\phi)}{1-\phi}-\frac{2\phi(1-\phi^T)}{(1-\phi)^2}.
  \label{eq:ar-sum}
\end{equation}
\end{proposition}

每个时期的最优病例量为 $c_0/c_1$ 与平衡规模的几何平均，与两阶段抽样中二级单元最优数量的形式相同\citep{cochran1977sampling}。$m$ 超过 $m^*$ 后，追加预算用于增加时期数比用于增加病例数更能降低误差。整数时期数、最低时期数和规模上限需另行约束，并与边界解比较。环境存在自相关时，一阶条件变为 $m^2=c_0A/(c_1v_Rg_T)$，其中 $g_T$ 由式~(\ref{eq:ar-sum}) 对 $T$ 求导得到；这一条件是必要条件，相应的预算阈值见附录~\ref{app:budget-ar}。

固定 $n$ 而改变 $\rho$ 的预算问题与上式不同。令 $H=v_R+(\Rbar^2+v_R)/(kn)$、$Q=\Rbar/n$，单期成本为 $c_0+c_\rho\rho$，则单期方差为 $H+Q/\rho$，内点最优报告概率为 $\rho^*=\sqrt{c_0Q/(c_\rho H)}$，再截取到 $(0,1]$。即使成本仍以每例报告病例 $c_1$ 计价，此路径也有 $c_\rho=c_1n$，不能套用式~(\ref{eq:optimal-m})（附录~\ref{app:prop3}）。

报告异质性可能逆转覆盖规模与报告完整度作为投入方向的优先次序。

\begin{corollary}[覆盖规模与报告完整度的比较]
\label{cor:switch}
若 $\tau$ 不依赖于 $\rho$，扩大覆盖规模 $n$ 对 $V_H$ 的降幅大于 $\rho$ 同比例提高时的降幅。若报告变异来自对 $\nu$ 个单位的二项抽样，即 $\Var(p_t)=\rho(1-\rho)/\nu$，则在下式成立时提高 $\rho$ 更为有利
\begin{equation}
  m>\frac{\rho^2(\nu-1)}{k},
  \label{eq:r3-switch}
\end{equation}
其中 $\nu$ 为每个时期抽样的报告单位数。若报告概率服从 Beta 先验分布 $p_t\sim\mathrm{Beta}(\rho\nu,(1-\rho)\nu)$，其方差为 $\rho(1-\rho)/(\nu+1)$，对应条件为 $m>\rho^2\nu/k$。
\end{corollary}

理论量的基线取值为 $\Rbar=1.15$、$v_R=0.04$、$k=0.35$、$\rho=0.25$ 与 $\tau=0.15$，闭式量及其适用条件列于表~\ref{tab:s7}。

\subsection{数据与分析窗口}
\label{sec:meth-data}

每周实验室确诊流感住院人数及报告医院百分比取自美国国家医疗安全网络（NHSN）的 Weekly Hospital Respiratory Data 数据表，覆盖美国 50 个州和哥伦比亚特区\citep{cdc2026nhsn}，所用版本于 2026-10-20 获取。数据表共含 2020-08-08 至 2026-10-10 的 21641 个辖区周。2024 年 4 月 30 日前为强制报告期，2024 年 5 月 1 日至 10 月 31 日为自愿报告期，2024 年 11 月 1 日起恢复强制报告\citep{cms2024reporting}；在三个报告期内，报告规模不低于五的辖区周分别为 6009、420 和 4077 个。

分析采用两组并列的早期增长窗口。规则窗口覆盖 2022--23 至 2025--26 四个流行季。每个窗口自 10 月 1 日起（2024 年自强制报告恢复后的 11 月 9 日起），包含全国周增长比不低于 1.05 的连续周，并于全国首个转折点的前一周结束，依次为 2022-10-01 至 2022-11-26、2023-10-07 至 2023-12-23、2024-11-09 至 2024-12-28 以及 2025-10-04 至 2025-12-27，共 42 个日历周。日历窗口覆盖 2022--23 至 2024--25 三个流行季的固定日期区间，依次为 2022-10-08 至 2022-11-26、2023-10-14 至 2023-12-16 以及 2024-11-09 至 2025-01-18，共 29 个日历周。两组窗口在 2022--23 与 2023--24 季大部分重叠，在 2024--25 季的终点不同，并分别报告。所有窗口均未纳入以零填补的病例数。观察性部分按 STROBE 声明的要素报告。

\subsection{周增长统计量}
\label{sec:meth-stats}

对每个州和每一周计算两个周增长统计量。加权增长比定义为
\begin{equation}
  G_t=\frac{C_{t+1}}{D_t},\qquad D_t=0.65\,C_t+0.25\,C_{t-1}+0.10\,C_{t-2}, % qa-allow
  \label{eq:dt}
\end{equation}
其中 $C_t$ 为第 $t$ 周的住院人数，$D_t$ 为加权报告规模，权重满足 $\sum_jw_j{}^2=0.495$。$G$ 仅使用 $D_t\ge5$ 的周。更新型比值将单周增长率通过代际间隔矩母函数代入 \citet{wallinga2007how} 式~(2.9) 换算为再生数闭式解，其定义为
\begin{equation}
  R_{A,t}=\left(1+r_t\theta\right)^{\alpha},\qquad r_t=\frac{1}{7}\log\frac{C_{t+1}+0.5}{C_t+0.5}, % qa-allow
  \label{eq:ra}
\end{equation}
其中 $r_t$ 为日增长率；代际间隔取均值 3.6 天、标准差 1.6 天的伽马分布（取 \citet{cowling2009estimation} 估计的家庭序列间隔），形状参数 $\alpha=(3.6/1.6)^2=5.0625$，尺度参数 $\theta=1.6^2/3.6=0.711111$ 天。$R_A$ 只在 $D_t\ge5$ 且 $C_t\ge5$ 的周计算。$R_A$ 由观测增长变换得到，并非 $R_t$ 的后验估计\citep{cori2013new}；两种统计量都属于表观量。

下文以 $M$ 泛指相对损失模型中的规模变量：$R_A$ 只含 $C_{t+1}$ 与 $C_t$，取 $M=C_t$；$G$ 的分母为三周加权和，取 $M=D_t$。两者都是住院代理，不等同于理论中的报告后有效感染压力 $m$。

\subsection{期望损失与表观平衡规模}
\label{sec:meth-loss}

常规监测中原发病例数无法观测，继发感染来源于感染压力 $\Lambda_{t+1}=\sum_{j\ge1}w_jI_{t+1-j}$\citep{fraser2007estimating, cori2013new}。给定历史 $\mathcal{F}_t$，对 $C_{t+1}$ 应用全方差公式，并令 $m=\rho\E(\Lambda_{t+1}\mid\mathcal{F}_t)$，相对预测方差可写为
\begin{equation}
  \frac{\Var(C_{t+1}\mid\mathcal{F}_t)}{m^2}=v_R+\frac{A}{m}+\mathrm{CV}^2(\Lambda_{t+1}\mid\mathcal{F}_t)\,(\Rbar^2+v_R),
  \label{eq:info-decomp}
\end{equation}
其中 $\mathrm{CV}$ 为变异系数。最后一项是另一个与规模无关的分量，仅在感染压力已知时消失。对于中心为 $\bar{X}$ 的统计量 $X_t$，相对损失 $L_t=[(X_t-\bar{X})/\bar{X}]^2$ 的模型设定为
\begin{equation}
  \E(L_t\mid M_t)=a+\frac{b}{M_t},
  \label{eq:affine}
\end{equation}
称 $b/a$ 为表观平衡规模，以区别于理论平衡规模 $\mx=A/v_R$。在一代恰为一周、无额外延迟、报告率固定已知且分母不含噪声的理想情形下，$a=v_R/\Rbar^2$、$b=A/\Rbar^2$，$b/a=\mx$。实际数据中 $D_t$ 经由感染到入院的延迟、周聚合和报告形成，分母本身带有报告误差，$R_A$ 还对增长作非线性变换，因此 $b/a$ 与 $\mx$ 之间需要显式校准，二者并不恒等（附录~\ref{app:latent-scale}）。若分母中只含二项噪声，其相对方差为 $(1-\rho)\sum_jw_j{}^2/m$，这一噪声使斜率近似增大为
\begin{equation}
  b\approx\frac{A+(1-\rho)(\Rbar^2+v_R)\sum_jw_j{}^2}{\Rbar^2}.
  \label{eq:b-delta}
\end{equation}
计算潜在平衡规模的情景范围时，$\Rbar$ 取 $G$ 的样本均值，$v_R\in\{0.04,0.08\}$，$k\in\{0.2,0.5\}$ 为情景取值（\citet{lloydsmith2005superspreading} 报告了 SARS 的 $k=0.16$），每例感染的住院概率取 $\rho\in\{0.0056,0.021\}$（由 \citet{tokars2018seasonal} 所报有症状感染至住院乘数之倒数确定）；这些情景下的平衡规模为 21.3--49.2 例。

中心 $\bar{X}$ 的定义构成以下层级。每一层在各自的样本上估计，并在每次自助重复中重新估计。
\begin{enumerate}[label=(\arabic*),leftmargin=2em]
  \item 主分析：各设定合并全部窗口的样本均值。
  \item 去除全国周效应：同一日历周内各州的均值，偏差乘以小样本因子 $n_w/(n_w-1)$，$n_w$ 为该周的州数。
  \item 州--季加性：以交替投影估计的周效应与州--季效应之和。
  \item 分季：各流行季窗口内的均值，用于分季截距。
  \item 报告完整度：未调整与覆盖校正的统计量在同一覆盖有效子样本上各按自身均值中心化，并使用同一规模 $M$。
\end{enumerate}

\subsection{估计与推断}
\label{sec:meth-inference}

式~(\ref{eq:affine}) 采用可行广义最小二乘拟合，权重为 $1/\hat{\mu}^2$，其中 $\hat{\mu}$ 为上一次迭代的拟合均值，迭代直至收敛\citep{wooldridge2010econometric}。同一州的观测在时间上相关，同一周的观测共享全国性影响，因此不确定性采用州 $\times$ 周双向自助法估计\citep{cameron2008bootstrap, cameron2011robust}。每次重复对 51 个州有放回抽样一次，对各流行季窗口内的周分别有放回抽样，以州抽取次数与周抽取次数之积作为观测权重。同一组抽取次数在 $R_A$ 与 $G$ 之间共享，中心在每次重复中重新估计。每个设定使用 2000 次重复，种子为 20261029，区间为百分位区间。为说明聚类结构的影响，另报告仅对州重抽样所得截距标准差与双向标准差之比。

截距分解依次以第~\ref{sec:meth-loss}~节层级中的（2）与（3）重新拟合式~(\ref{eq:affine})。全国周份额定义为 $(a-a_{\text{周}})/a$，州--季分量为加入州--季效应后截距的进一步变化。作为对函数形式的描述，同时拟合无截距幂律 $\E(L_t\mid M_t)=cM_t{}^{-\gamma}$。仿射模型的局部对数斜率为 $-1/[1+M/(b/a)]$，其在观测规模上的平均值是几何描述；$\gamma$ 由另一全局拟合目标决定，与平均局部斜率之间不存在恒等关系。幂律拟合在部分自助重复中不收敛到内点，自助分布因此出现重尾，以标准差为基础的区间与百分位区间不一致。规则窗口报告 $\gamma$ 的百分位区间及不收敛重复的比例；日历窗口不收敛比例超过 13\%，只报告点估计。本文不对 $\gamma=1$ 作检验，截距是否为正与两种拟合在分箱均值上的接近程度一并作为描述性证据。

\subsection{报告完整度}
\label{sec:meth-reporting}

将各州各周的住院人数除以对应的报告医院比例，得到覆盖校正统计量。检验对象为 $\Delta a=a_{\mathrm{raw}}-a_{\mathrm{adj}}$，即同一覆盖有效子样本、同一规模 $M$ 上两个按自身均值中心化的 FGLS 截距之差。等效性采用 5\% 水平的双单侧检验（TOST）：$\Delta a$ 的 90\% 配对双向自助区间完全落在 $(-\delta,\delta)$ 之内时判为等效\citep{schuirmann1987comparison}。

等效界 $\delta$ 须与截距统计量处于同一尺度，因此由配对校准得到，而不直接取命题~\ref{prop:report-noise} 中的 $v_H$。每次校准重复中，对每个州周抽取 $U_{it}\sim\mathrm{Uniform}(-\sqrt3,\sqrt3)$，把观测住院数乘以 $(1+\tau_0U_{it})$。同一冲击同时进入它所在周作为分子和作为分母的各项，规模 $M$ 保持不变。随后按与实证分析相同的变换、中心化与 FGLS 重新拟合截距。$\delta$ 取 2000 次重复中截距增量的平均值，$\tau_0=0.10$ 与 0.05 两个水平分别计算，并报告其蒙特卡罗标准误。这一界因此对应独立有界乘性报告冲击模型下的敏感性问题。它检验观测到的截距变化是否小于相对波动为 $\tau_0$ 的独立报告冲击所能产生的截距变化。合并样本的检验为主要结果，分季检验作为描述。医院报告比例是机构参与度的代理，不包括报告医院内部检测强度与编码的变化。

自愿报告期只有 2024 年 5 月至 10 月一个非流行季，样本中低规模州周占多数，FGLS 截距在这一时段不稳定（自愿报告期的点估计为负）。对这一对照改用稳健的离散度量：比较 2023 年、2024 年与 2025 年 5 月至 10 月间 $\log[(C_{t+1}+0.5)/(D_t+0.5)]$ 的四分位距。区间来自 2000 次双向自助重复（州抽取次数在三年间共享，周在年内抽样，种子 20261941）。为排除规模构成差异，另按 2023 年与 2025 年合并样本中 $\log D_t$ 的十分位分布对各年重新加权，并按日历月逐月比较。

\subsection{流行季与分母权重}
\label{sec:meth-seasons}

各流行季截距采用 DerSimonian--Laird 随机效应模型合并\citep{dersimonian1986meta}；合并截距采用自由度为 $S-1$ 的 $t$ 区间，季间方差采用 Q-profile 区间\citep{viechtbauer2007confidence}。推论~\ref{cor:season} 给出任一季间方差估计相对标准误的下界\citep{searle1992variance}。截距比 $q_1=a_{(1,0,0)}/a_{w}$ 与 $q_2=a_{(1/3,1/3,1/3)}/a_{w}$ 分别将单周分母与等权三周分母同式~(\ref{eq:dt}) 的权重相比较，在规则窗口中计算，区间来自与 $G$ 相同的双向自助抽取。比较基准为纯报告噪声零模型：以平滑后的全国流行季曲线和各州固定份额给出期望住院数，乘以对数正态报告率（$\mathrm{CV}^2\in\{0.0036,0.0143,0.0287,0.0574\}$）后作泊松抽样，不含环境变异。在每条模拟的全国路径上按同一规则重新选择窗口，每个噪声水平重复 300 次。零模型范围取各噪声水平下比值均值的极差，观测比值落在范围之外的双侧 $p$ 值由自助分布计算。

\subsection{通用独立面板模拟}
\label{sec:meth-sim}

理论结果与表观统计量的性质通过两类蒙特卡罗模拟检验。图~\ref{fig:theory} 中的每个点基于 20000 次重复。通用独立面板设计中，表观周统计量由日尺度更新模型模拟生成，取 $\Rbar=1.2$，代际间隔同上。每周 $R_t$ 在各周间独立地从方差为 $v_R$ 的伽马分布中抽取；后代服从负二项分布；住院经均值 5.5 天、标准差 3.0 天的伽马延迟后以概率 $\rho_h$ 二项发生，延迟分布的参数为假定值。每个重复包含 500 个州，初始报告规模在 2 与 150 之间服从对数均匀分布，另含四个预热周与六个观测周。十七个情景由 $v_R\in\{0.01,0.02,0.04,0.08,0.16\}$ 与 $k\in\{0.1,0.35,1\}$ 在 $\rho_h=0.01$ 下交叉组合，并在基线设定下增加 $\rho_h\in\{0.005,0.02\}$，每个情景重复 20 次。每个截距的无噪声极限来自 20000 个州的平均场运行，其中全部计数噪声均以期望值代替。

\subsection{潜在环境方差的模拟反演}
\label{sec:meth-inversion}

作为补充分析，实证截距经两步映射为潜在环境方差 $v_R$。第一步，以数据匹配的模拟器给出各统计量期望拟合截距随 $v_R$ 变化的函数。每个模拟面板的横截面与长度和分析面板相同，包含 51 个州和四个规则窗口，观测周数依次为 9、12、8 和 13。流行季均值 $R_s$ 取各窗口全国周增长几何均值的更新型变换，各州初始住院数等于该州在窗口首周的观测住院数。周再生数为 $R_{i,s,w}=R_s+\eta_{s,w}+\varepsilon_{i,s,w}$，其中全国周冲击 $\eta$ 的方差为 $0.45\,v_R$，为所有州共有；州冲击 $\varepsilon$ 的方差为 $0.55\,v_R$；取值在 0.05 处截断。0.45 取在两种统计量观测到的全国周份额之间，因此模拟对组内相关系数与全国周份额的复现只是一致性检查。传播为后代服从负二项分布（$k=0.35$）的日尺度更新过程，住院经上述延迟后以概率 $\rho_h=0.01$ 二项发生；两种统计量的计算与筛选完全按数据处理，分析者采用均值 3.6 天的代际间隔。设计取 $v_R\in\{0.02,0.04,0.06,0.09\}$，每个取值 200 个面板，对每个面板拟合式~(\ref{eq:affine})。校准线 $\E(\hat{a})=\gamma_0+\gamma_1v_R$ 采用加权最小二乘拟合，权重为各均值蒙特卡罗方差的倒数；尺度因子 $\kappa=\E(\hat{a})/(v_R/\Rbar^2)$ 中的 $\Rbar$ 取按观测行加权的真实均值。

第二步，对各统计量计算隐含方差 $v_R^{\mathrm{impl}}=(a-\gamma_0)/\gamma_1$。其区间合并两部分变异：两个实证截距的 2000 次联合双向自助抽样（两者共用州与周的抽取次数，种子 20261921），以及 $(\gamma_0,\gamma_1)$ 按其蒙特卡罗抽样分布的独立抽取（种子 20261922）。两种统计量的一致定义为两者之差的 95\% 区间包含零，即未拒绝。敏感性设定每次只改变设计中的一个要素（表~\ref{tab:s8}）。全国冲击被去除或份额提高到 0.65，周间一阶自相关取 0.5 或 0.25，真实代际间隔均值取 2.2、2.9 或 4.3 天而分析者仍用 3.6 天。其余设定把延迟均值改为 3.5 天或 8 天，把 $k$ 降至 0.1，以作用于报告计数的同步乘性冲击代替传播冲击，或令每周住院数取自方差均值比为 1.5 或 2.0 的伽马--泊松混合（表~\ref{tab:s9}）。周冲击的独立性以各统计量 $x/\bar{x}-1$ 在州--季节块内的汇总滞后一阶自相关检验，区间来自 1000 次按州整群的自助抽样，并与 $v_R=0.04$ 时模拟面板中的分布比较。校准后的理论 $b/a$ 采用同一设计，取 $v_R\in\{0.04,0.08\}$、$k\in\{0.2,0.5\}$ 与 $\rho_h\in\{0.0056,0.021\}$；单个面板间 $b/a$ 的取值范围给出同等规模单个面板的预期分布。模拟种子由 20261020 按设定、$v_R$ 取值与重复块索引派生。
